% Lecture visuelle directe avec la vision du modèle ; aucun OCR.
% Source : PDF original, vue 001, côté droit, page manuscrite 1.
% Le côté gauche de la vue 001 ne porte que des transparences de l'autre face.
% Statut : transcrit ; relecture indépendante à effectuer.
\chapter*{Épître dédicatoire à la Sainte Vierge}
\addcontentsline{toc}{chapter}{Épître dédicatoire à la Sainte Vierge}
\markboth{Épître dédicatoire}{Épître dédicatoire}

\begin{center}
\textit{À Marie !}\par
À ma bonne mère ! Vierge immaculée ! Auguste protectrice et patronne de la Congrégation de Lyon.
\end{center}

Ô ma mère ! Vous qui, dans les peines et les vicissitudes de la vie, avez toujours été mon refuge et mon soutien, permettez à votre indigne serviteur de vous offrir et de vous consacrer ce travail, qu'il entreprend aujourd'hui pour l'édification de ses frères.

C'est avec une confiance toute filiale que j'ose m'adresser à vous et vous dire dans toute la simplicité de mon cœur : Souvenez-vous, ô ma mère, que vous m'avez adopté, il y a plus de 43 ans, dans le sanctuaire érigé en votre honneur à Notre-Dame des Hermites ; et que dès lors je suis devenu votre enfant. Souvenez-vous qu'au 3\up{e} dimanche de juillet 1802, vous avez daigné renouveller cette adoption d'une manière plus solennelle encore.

J'ai donc acquis le droit, ô ma mère, de solliciter une petite place dans votre cœur maternel. C'est là que je puis espérer de trouver ces sentimens d'amour et de reconnaissance qui, reproduits sur le papier, parviendront aux cœurs de mes frères. Daignez, je vous en conjure, ô Marie, m'obtenir les grâces dont j'ai besoin, pour rendre l'ouvrage que je commence vraiment
% Source : PDF original, vue 002, côté gauche, page manuscrite 2.
utile à tous les Congréganistes. Guidez ma plume ; que par vous elle retrace et conserve le souvenir de toutes les merveilles qui ont été opérées en notre faveur, et de toutes les grâces qui ont été répandues, par votre intercession, sur la Congrégation et sur tous ceux qui ont été assez heureux pour en faire partie.

Ce n'est qu'en me plaçant ainsi sous vos auspices, et en espérant votre puissant secours, que j'ose tenter d'élever un monument à la gloire de votre divin fils, ainsi qu'à la vôtre, pour rappeler d'âge en âge la mémoire de tant de bienfaits et exciter de plus en plus tous les cœurs à la reconnaissance. Que nos plus jeunes frères, et que tous ceux qui dans la suite viendront prendre place dans notre heureuse famille, apprennent par quelles grâces toutes de prédilection leur bonheur a été préparé, et qu'avec leurs aînés ils bénissent le Seigneur et glorifient son saint nom.

\newpage
Daignez donc, ô ma mère, agréer cet hommage. Que par vous il parvienne au cœur de votre divin fils. Qu'il daigne répandre ses bénédictions sur l'ouvrage et sur le misérable instrument dont il veut bien se servir pour le rédiger. Que par vous, ô Marie, ce pauvre serviteur de Dieu, votre enfant, obtienne le pardon de toutes les fautes dont il a pu se rendre coupable, une soumission entière à la volonté de son souverain maître et un désir sincère d'expier par la pénitence tout ce qui dans sa vie entière a pu lui déplaire. Obtenez-lui surtout la grâce de mourir dans l'amour de Jésus. Venez, refuge des pécheurs, venez à son secours dans ce moment suprême, afin qu'abordant heureusement au port de l'éternité, il puisse jouir à jamais de la vie de la charité et glorifier avec les anges et les saints et le fils et la mère, pendant
% Source : PDF original, vue 002, côté droit, page manuscrite 3, haut de page.
les siècles des siècles. Ainsi soit-il !

\begin{center}
Vive Jésus ! Vive Marie !\par
\textit{Monstra te esse Matrem.}\footnote{Hymne \textit{Ave maris stella}, quatrième strophe : « Montrez que vous êtes mère. » (Note du transcripteur).}
\end{center}
